% Generated draft: requires mathematical and source review.
\documentclass[11pt]{article}
\usepackage[margin=25mm]{geometry}
\usepackage{fontspec}
\setmainfont{Latin Modern Roman}
\usepackage{amsmath,amssymb,mathtools,unicode-math}
\setmathfont{Latin Modern Math}
\usepackage{enumitem,needspace,xurl,hyperref}
\setlength{\parindent}{0pt}
\setlength{\parskip}{5pt}
\setlength{\emergencystretch}{3em}
\newcommand{\N}{\mathbb{N}}
\newcommand{\Z}{\mathbb{Z}}
\newcommand{\Q}{\mathbb{Q}}
\newcommand{\R}{\mathbb{R}}
\newcommand{\C}{\mathbb{C}}
\newcommand{\F}{\mathbb{F}}
\newcommand{\A}{\mathbb{A}}
\newcommand{\PP}{\mathbb{P}}
\newcommand{\E}{\mathbb{E}}
\newcommand{\eps}{\varepsilon}
\newcommand{\dd}{\,\mathrm{d}}
\begin{document}
\section*{UnsolvedMath ID: 30005750}
\textbf{Source problem code:} catalogue entry
\subsection{Definitions and mathematical statement}
The supplied record poses a classification problem rather than a single fully specified proposition. A decomposition of continuity is an equivalence of the form
\[
 f\text{ is continuous}\quad\Longleftrightarrow\quad W_1(f)\land\cdots\land W_m(f),\qquad m\geq 2,
\]
where the predicates \(W_i\) express weaker continuity properties on the relevant function space. The source describes many such decompositions, historically going back to Baire, but supplies neither a list of the predicates nor their definitions.

For a real-valued function \(f:[0,1]\to\R\) and real numbers \(0\leq p\leq q\leq1\), put
\[
 A_{f,p,q}:=\{f(y):y\in[p,q]\}.
\]
When this set is bounded above, its real supremum is characterized by
\[
 s=\sup A_{f,p,q}
 \quad\Longleftrightarrow\quad
 \bigl(\forall y\in[p,q]\;f(y)\leq s\bigr)
 \land
 \bigl(\forall\varepsilon>0\;\exists y\in[p,q]\;f(y)>s-\varepsilon\bigr).
\]
The case \(p=q\) gives \(s=f(p)\). If \(A_{f,p,q}\) is unbounded above, it has no real supremum; allowing the value \(+\infty\) would require an explicitly extended-real formulation.

\textbf{Problem scheme.} For each weak continuity predicate \(W\) occurring in the intended decompositions of ordinary continuity, determine the computational or logical strength of obtaining the interval-supremum functional
\[
 S_W(p,q,f)=\sup_{y\in[p,q]}f(y),\qquad W(f).
\]
The displayed real-valued specification applies when the range on the interval is bounded above. Whether boundedness is a hypothesis, follows from the particular predicate, or is replaced by an extended-real output is unresolved in the supplied excerpt. The source also does not specify whether the endpoints must be rational, how functions and real outputs are represented, or which precise notion of computational difficulty is intended.

The excerpt places the question in higher-order reverse mathematics and reports that, for most of the weak continuity notions under discussion, the restricted supremum functional can be defined in \(\mathsf{RCA}_0^{\omega}+(\exists^2)\). Here \((\exists^2)\) conventionally denotes the existence of a functional \(\Phi:\N^{\N}\to\N\) satisfying
\[
 \forall g:\N\to\N\quad
 \bigl(\Phi(g)=0\ \Longleftrightarrow\ \exists n\in\N\;g(n)=0\bigr).
\]
This is a supplied contextual claim; the excerpt does not identify its precise class of predicates or the formal meaning of definability being used.

The excerpt additionally asserts the existence of a function \(f:[0,1]\to\R\) that is not Baire class \(2\), discusses codes for Baire class \(3\) functions, and reports a consistency claim that all functions are Baire class \(2\) in a theory rendered as \(\texttt{Z}\omega\texttt{2}\). The ambient theories, models, coding conventions, and exact notation needed to interpret these remarks are not supplied. These contextual assertions are retained as source qualifications and are not adopted here as unqualified mathematical premises.
\subsection{Short English statement}
Classify how difficult it is to obtain suprema of real-valued functions satisfying the weak continuity properties appearing in decompositions of continuity. A precise classification requires the missing definitions, representations, endpoint conventions, and treatment of unbounded functions. The supplied record labels the problem partially solved; its dated assessment says that many supremum principles have been classified but the broader taxonomy remains incomplete.
\subsection{Sources}
\begin{itemize}
\item Local UnsolvedMath registry: supplied complete source record, permanent ID \(30005750\), title \emph{Supremum Computation for Weakly Continuous Functions}, exact source problem code \texttt{catalogue entry}. The record belongs to the Oberwolfach Reports Open Problems collection. No individual proposer is supplied; the proposed year is recorded as 2023.
\item Supplied report citation: \emph{Mathematical Logic: Proof Theory, Constructive Mathematics} (2024), Oberwolfach Reports, pp.\ 3013--3061, DOI \texttt{10.4171/owr/2023/53}. Original OWR report: \url{https://doi.org/10.4171/owr/2023/53}. The background describes these reports as documenting mathematical workshops and their open-problem discussions.
\item Supplied literature reference, labelled ``weak-continuity paper'': \url{https://arxiv.org/abs/2405.06420}. The supplied record gives no authors, title, or theorem identifiers for this reference.
\item Literature assessment checked at \texttt{2026-08-22}: supplied dated curation, not independently verified. The supplied assessment states that a primary reverse-mathematics paper classifies supremum principles for many weak-continuity notions, while the broader taxonomy question in the excerpt is not exhausted. The supplied status is \texttt{partially\_solved}; the background explicitly requests independent verification of time-sensitive status and cited claims.
\item The registry labels the statement \texttt{corrected\_verified} and reports that it was checked against the official Oberwolfach report. This is supplied verification metadata, not an independent verification performed for this draft.
\end{itemize}
\end{document}
